\scanpage{69}
Donc $T$ est connue (\unsure{à|pr} équiv. près) quand on connaît $S$ et les $r_i:S\longrightarrow(\mathrm{Ens})$ [et on connaît alors \unsure{même|$=$} $T$ sur $\tau^{I}$ \ldots], ou \deleted{de} façon plus intrinsèque, si on connaît \deleted{\unsure{$R\subset$}} \inserted{l'image de $R$ dans} $\underline{\mathrm{Hom}}(S,\mathrm{Ens})=\widehat{S^\circ}$.

\textcircled{3°} On a
\[
S\overset{\sim}{\simeq}\underline{\mathrm{Hom}}_{\Delta}(R,(\mathrm{Ens}))\subset\widehat{R^\circ},
\]
et cette sous-catégorie \uline{contient} $R^\circ$ [car les foncteurs \inserted{cov.} représentables \deleted{\unsure{contra}} de $R$ dans $\mathrm{Ens}$ commutent aux $\varprojlim$ !) donc on a
\[
\boxed{R^\circ\overset{\unsure{j|\partial|\partial\,\xi}}{\subset}S}
\]
et le foncteur \inserted{$\widetilde r:S\longrightarrow(\mathrm{Ens})$} associé à un $r\in\mathrm{Ob}\,R$ est le foncteur qu'il représente \uline{dans} $S$. Ainsi $R$ est connue quand on connaît son image \unsure{ess.|ens.} dans $\underline{\mathrm{Hom}}(S,\mathrm{Ens}]=\widehat{S^\circ}$, car $R\longrightarrow\deleted{\underline{\mathrm{Hom}}}\widehat{S^\circ}$ est le composé $R\xrightarrow{j^\circ}S^\circ\longrightarrow\widehat{S^\circ}$ de deux fonctrs pl. fid.

Ainsi $S=T_{\deleted{\unsure{R}}}(\mathrm{Ens})$ est muni d'une sous-catégorie str. pleine $\Sigma$, ($\approx R^\circ$) [formée des $\varphi\in\mathrm{Ob}\,S$ tels que \deleted{pour} $\forall C\in\mathrm{Ob}\,\Delta$ et tout $\xi\in T(C)$, définissant $\widetilde\xi:C^\circ\longrightarrow S$, \deleted{les} le composé $\inserted{x\mapsto}\mathrm{Hom}(\varphi,\overline{\xi}(x))$ de $C^\circ$ dans $\mathrm{Ens}$ soit représentable]. \inserted{La sous-catégorie} $\Sigma$ est stable dans $S$ par \margin{d'où un foncteur $T(C)\xrightarrow{\unsure{\widetilde\varphi_C|\varphi_C}}C$} $\varinjlim$ (\uline{inductives}) de type $d'$, \deleted{\unsure{de|et} pour tout} On trouve ainsi une équivalence \inserted{explicite} entre $\Sigma^\circ$ et $R$.

\deleted{\textcircled{4°} Soit \deleted{\unsure{$d$}} $\Theta'$ l'\unsure{ens|un|une} de tous les petits diagrammes, $\Delta_0$ \deleted{\unsure{l'une}} la \inserted{2-}catégorie \unsure{de|des} toutes les catégories où les $\varprojlim$ de type $\Theta'$ existent, Supposons que dans la définition de $\Delta$ ne figurent pas des conditions de finitude, de sorte que $\deleted{\unsure{\Delta}}\Delta_0\subset\Delta$. Donc la théorie $T$ \inserted{$=T_R$} 2-repr pour $\Delta$ définit une théorie pour $\Delta_0$, dont nous allons montrer la 2-représentabilité. \deleted{Soit $C\in\mathrm{Ob}\,\Delta_0$, on a} \uline{Si} c'est représentable, on}
\[
\deleted{\underline{\mathrm{Hom}}_{\Delta_{\unsure{0|4}}}(R,C)\simeq\underline{\mathrm{Hom}}_{\varinjlim}(R^\circ,C^\circ)^\circ\unsure{\simeq|\sim}}
\]
\deleted{sait que la catégorie $R_0$ qui représente \deleted{\unsure{sera|dans}} une catégorie avec \deleted{$\varinjlim$} $\varprojlim$ quelconques, \deleted{\unsure{muni de}} etc}
% archive sheet number: 68
